% GENERATED FILE. Edit canonical lessons, then run npm run generate:books.
% canonical-source-hash: fnv1a64:6956492d494d76a0
% canonical-lessons: TA-C520-ceyarkaikkol, TA-C520-taniyanki, TA-C520-calavai-iyantiram, TA-C520-mincatanam, TA-C520-istiri-potu

\chapter{Machines in the Home}
\label{ch:ta-machines-in-the-home}
\hlchaptermodality{520}

\begin{chapteropening}
\textbf{By the end of this chapter:} \emph{I can talk about machines that work for us.}

Complete the last lesson of chapter 520: I can talk about machines that work for us.
\end{chapteropening}

\section[\texorpdfstring{\ta{செயற்கைக்கோள்} (ceyaṟkaikkōḷ)}{செயற்கைக்கோள் (ceyaṟkaikkōḷ)}]{\ta{செயற்கைக்கோள்} (ceyaṟkaikkōḷ) — a satellite}

\label{lesson:TA-C520-ceyarkaikkol}



\begin{quote}
Before the new one: say the Tamil for a printer, then the Tamil for electronic.
\end{quote}

\subsection*{You'll want to know: \ta{செயற்கைக்கோள்}}
\textbf{\ta{செயற்கைக்கோள்}} — \emph{ceyaṟkaikkōḷ} — \textquotedblleft{}a satellite\textquotedblright{}.

\begin{quote}
\ta{செயற்கைக்கோள்} \ta{மூலம்} \ta{வானிலைப்} \ta{படங்கள்} \ta{கிடைக்கின்றன}.
\end{quote}

\begin{quote}
\emph{ceyaṟkaikkōḷ mūlam vāṉilaip paṭaṅkaḷ kiṭaikkiṉṟaṉa.}
\end{quote}

\begin{quote}
Weather pictures come from a satellite.
\end{quote}

\ta{செயற்கை} \textquotedblleft{}man-made\textquotedblright{} + \ta{கோள்} \textquotedblleft{}a planet\textquotedblright{}: a man-made planet.

\begin{cousinweb}
\textbf{\ta{செயற்கைக்கோள்}} (\emph{ceyaṟkaikkōḷ}) is a modern compound: \ta{செயற்கை} (\emph{ceyaṟkai}), \textquotedblleft{}man-made\textquotedblright{}, from \ta{செய்} (DEDR 1957), and \ta{கோள்} (\emph{kōḷ}), \textquotedblleft{}a planet\textquotedblright{}; the origin of \ta{கோள்} is not settled.
\end{cousinweb}

\begin{grammarlens}[title={What you've built}]
One more word for describing what you lived and felt.
\end{grammarlens}

\subsection*{Your turn}
\begin{itemize}
\raggedright
  \item \emph{Say it:} \emph{ceyaṟkaikkōḷ}
  \item \emph{Say it:} \emph{ceyaṟkaikkōḷ}, once more
  \item \emph{Say it:} \emph{ceyaṟkaikkōḷ mūlam vāṉilaip paṭaṅkaḷ kiṭaikkiṉṟaṉa.} — Weather pictures come from a satellite.
  \item \emph{Recall:} say the Tamil for a printer, then the Tamil for electronic, then say \emph{ceyaṟkaikkōḷ} again
\end{itemize}

\begin{tcolorbox}[breakable,colback=teal!4,colframe=teal!35!black,title={Before you move on}]
What does \ta{செயற்கைக்கோள்} mean? (\textquotedblleft{}A satellite\textquotedblright{}.) Say it once more.
\end{tcolorbox}
\section[\texorpdfstring{\ta{தானியங்கி} (tāṉiyaṅki)}{தானியங்கி (tāṉiyaṅki)}]{\ta{தானியங்கி} (tāṉiyaṅki) — automatic; an automatic machine}

\label{lesson:TA-C520-taniyanki}



\begin{quote}
Before the new one: say the Tamil for electronic, then the Tamil for a satellite.
\end{quote}

\subsection*{You'll want to know: \ta{தானியங்கி}}
\textbf{\ta{தானியங்கி}} — \emph{tāṉiyaṅki} — \textquotedblleft{}automatic; an automatic machine\textquotedblright{}.

\begin{quote}
\ta{வங்கியின்} \ta{தானியங்கி} \ta{இயந்திரத்தில்} \ta{பணம்} \ta{எடுத்தேன்}.
\end{quote}

\begin{quote}
\emph{vaṅkiyiṉ tāṉiyaṅki iyantirattil paṇam eṭuttēṉ.}
\end{quote}

\begin{quote}
I took money out of the bank's automatic machine.
\end{quote}

It stands before a noun: \ta{தானியங்கி} \ta{இயந்திரம்}, an automatic machine; \ta{தானியங்கிக்} \ta{கதவு}, an automatic door.

\begin{cousinweb}
\textbf{\ta{தானியங்கி}} (\emph{tāṉiyaṅki}) is a modern native compound: \ta{தான்} (\emph{tāṉ}), \textquotedblleft{}oneself, itself\textquotedblright{} (DEDR 3196), and \ta{இயங்கி} (\emph{iyaṅki}), \textquotedblleft{}that which moves\textquotedblright{}, from \ta{இயங்கு}, \textquotedblleft{}to move\textquotedblright{} (DEDR 469).
\end{cousinweb}

\begin{grammarlens}[title={What you've built}]
One more word for describing what you lived and felt.
\end{grammarlens}

\subsection*{Your turn}
\begin{itemize}
\raggedright
  \item \emph{Say it:} \emph{tāṉiyaṅki}
  \item \emph{Say it:} \emph{tāṉiyaṅki}, once more
  \item \emph{Say it:} \emph{vaṅkiyiṉ tāṉiyaṅki iyantirattil paṇam eṭuttēṉ.} — I took money out of the bank's automatic machine.
  \item \emph{Recall:} say the Tamil for electronic, then the Tamil for a satellite, then say \emph{tāṉiyaṅki} again
\end{itemize}

\begin{tcolorbox}[breakable,colback=teal!4,colframe=teal!35!black,title={Before you move on}]
What does \ta{தானியங்கி} mean? (\textquotedblleft{}Automatic; an automatic machine\textquotedblright{}.) Say it once more.
\end{tcolorbox}
\section[\texorpdfstring{\ta{சலவை} \ta{இயந்திரம்} (calavai iyantiram)}{சலவை இயந்திரம் (calavai iyantiram)}]{\ta{சலவை} \ta{இயந்திரம்} (calavai iyantiram) — a washing machine}

\label{lesson:TA-C520-calavai-iyantiram}



\begin{quote}
Before the new one: say the Tamil for a satellite, then the Tamil for automatic.
\end{quote}

\subsection*{You'll want to know: \ta{சலவை} \ta{இயந்திரம்}}
\textbf{\ta{சலவை} \ta{இயந்திரம்}} — \emph{calavai iyantiram} — \textquotedblleft{}a washing machine\textquotedblright{}.

\begin{quote}
\ta{சலவை} \ta{இயந்திரத்தில்} \ta{துணிகளைப்} \ta{போட்டேன்}.
\end{quote}

\begin{quote}
\emph{calavai iyantirattil tuṇikaḷaip pōṭṭēṉ.}
\end{quote}

\begin{quote}
I put the clothes in the washing machine.
\end{quote}

\ta{சலவை} \textquotedblleft{}washing of cloth\textquotedblright{} + \ta{இயந்திரம்} (\emph{iyantiram}). The noun takes -\ta{அத்து}-: \ta{இயந்திரத்தில்}.

\begin{cousinweb}
\textbf{\ta{சலவை} \ta{இயந்திரம்}} (\emph{calavai iyantiram}) joins the native \ta{சலவை} (\emph{calavai}), \textquotedblleft{}washing, bleaching of cloth\textquotedblright{} (DEDR 2368), and \ta{இயந்திரம்} (\emph{iyantiram}), \textquotedblleft{}a machine\textquotedblright{}, from Sanskrit \emph{yantra} (CDIAL 10412).
\end{cousinweb}

\begin{grammarlens}[title={What you've built}]
One more word for describing what you lived and felt.
\end{grammarlens}

\subsection*{Your turn}
\begin{itemize}
\raggedright
  \item \emph{Say it:} \emph{calavai iyantiram}
  \item \emph{Say it:} \emph{calavai iyantiram}, once more
  \item \emph{Say it:} \emph{calavai iyantirattil tuṇikaḷaip pōṭṭēṉ.} — I put the clothes in the washing machine.
  \item \emph{Recall:} say the Tamil for a satellite, then the Tamil for automatic, then say \emph{calavai iyantiram} again
\end{itemize}

\begin{tcolorbox}[breakable,colback=teal!4,colframe=teal!35!black,title={Before you move on}]
What does \ta{சலவை} \ta{இயந்திரம்} mean? (\textquotedblleft{}A washing machine\textquotedblright{}.) Say it once more.
\end{tcolorbox}
\section[\texorpdfstring{\ta{மின்சாதனம்} (miṉcātaṉam)}{மின்சாதனம் (miṉcātaṉam)}]{\ta{மின்சாதனம்} (miṉcātaṉam) — an electrical appliance}

\label{lesson:TA-C520-mincatanam}



\begin{quote}
Before the new one: say the Tamil for automatic, then the Tamil for a washing machine.
\end{quote}

\subsection*{You'll want to know: \ta{மின்சாதனம்}}
\textbf{\ta{மின்சாதனம்}} — \emph{miṉcātaṉam} — \textquotedblleft{}an electrical appliance\textquotedblright{}.

\begin{quote}
\ta{புயலின்போது} \ta{எல்லா} \ta{மின்சாதனங்களையும்} \ta{அணைத்தோம்}.
\end{quote}

\begin{quote}
\emph{puyaliṉpōtu ellā miṉcātaṉaṅkaḷaiyum aṇaittōm.}
\end{quote}

\begin{quote}
During the storm we switched off all the electrical appliances.
\end{quote}

Plural: \ta{மின்சாதனங்கள்}, \emph{miṉcātaṉaṅkaḷ}.

\ta{மின்} \textquotedblleft{}electric\textquotedblright{} + \ta{சாதனம்} \textquotedblleft{}a device\textquotedblright{}. Plural: \ta{மின்சாதனங்கள்}.

\begin{cousinweb}
\textbf{\ta{மின்சாதனம்}} (\emph{miṉcātaṉam}) joins the native \ta{மின்} (\emph{miṉ}), \textquotedblleft{}electric\textquotedblright{} (DEDR 4876), and \ta{சாதனம்} (\emph{cātaṉam}), \textquotedblleft{}a means, a device\textquotedblright{}, from Sanskrit \emph{sādhana} (CDIAL 13334).
\end{cousinweb}

\begin{grammarlens}[title={What you've built}]
One more word for describing what you lived and felt.
\end{grammarlens}

\subsection*{Your turn}
\begin{itemize}
\raggedright
  \item \emph{Say it:} \emph{miṉcātaṉam}
  \item \emph{Say it:} \emph{miṉcātaṉam}, once more
  \item \emph{Say it:} \emph{puyaliṉpōtu ellā miṉcātaṉaṅkaḷaiyum aṇaittōm.} — During the storm we switched off all the electrical appliances.
  \item \emph{Recall:} say the Tamil for automatic, then the Tamil for a washing machine, then say \emph{miṉcātaṉam} again
\end{itemize}

\begin{tcolorbox}[breakable,colback=teal!4,colframe=teal!35!black,title={Before you move on}]
What does \ta{மின்சாதனம்} mean? (\textquotedblleft{}An electrical appliance\textquotedblright{}.) Say it once more.
\end{tcolorbox}
\section[\texorpdfstring{\ta{இஸ்திரி} \ta{போடு} (istiri pōṭu)}{இஸ்திரி போடு (istiri pōṭu)}]{\ta{இஸ்திரி} \ta{போடு} (istiri pōṭu) — to iron (clothes)}

\label{lesson:TA-C520-istiri-potu}



\begin{quote}
Before the new one: say the Tamil for a washing machine, then the Tamil for an electrical appliance.
\end{quote}

\subsection*{You'll want to know: \ta{இஸ்திரி} \ta{போடு}}
\textbf{\ta{இஸ்திரி} \ta{போடு}} — \emph{istiri pōṭu} — \textquotedblleft{}to iron (clothes)\textquotedblright{}.

\begin{quote}
\ta{அம்மா} \ta{என்} \ta{சட்டையை} \ta{இஸ்திரி} \ta{போட்டாள்}.
\end{quote}

\begin{quote}
\emph{ammā eṉ caṭṭaiyai istiri pōṭṭāḷ.}
\end{quote}

\begin{quote}
Mother ironed my shirt.
\end{quote}

Past (I …): \ta{இஸ்திரி} \ta{போட்டேன்}, \emph{istiri pōṭṭēṉ}.

The garment takes -\ta{ஐ}: \ta{சட்டையை} \ta{இஸ்திரி} \ta{போடு}. The iron itself is an \ta{இஸ்திரிப்} \ta{பெட்டி}.

\begin{cousinweb}
In \textbf{\ta{இஸ்திரி} \ta{போடு}} (\emph{istiri pōṭu}) the verb is the native \ta{போடு} (\emph{pōṭu}), \textquotedblleft{}to put\textquotedblright{}; the origin of \ta{இஸ்திரி} (\emph{istiri}), \textquotedblleft{}ironing\textquotedblright{}, is not settled.
\end{cousinweb}

\begin{grammarlens}[title={What you've built}]
One more word for describing what you lived and felt.
\end{grammarlens}

\subsection*{Your turn}
\begin{itemize}
\raggedright
  \item \emph{Say it:} \emph{istiri pōṭu}
  \item \emph{Say it:} \emph{istiri pōṭu}, once more
  \item \emph{Say it:} \emph{ammā eṉ caṭṭaiyai istiri pōṭṭāḷ.} — Mother ironed my shirt.
  \item \emph{Recall:} say the Tamil for a washing machine, then the Tamil for an electrical appliance, then say \emph{istiri pōṭu} again
\end{itemize}

\begin{tcolorbox}[breakable,colback=teal!4,colframe=teal!35!black,title={Before you move on}]
What does \ta{இஸ்திரி} \ta{போடு} mean? (\textquotedblleft{}To iron (clothes)\textquotedblright{}.) Say it once more.
\end{tcolorbox}
